
Our mechanistic attribution approach builds on three research threads: LLM-guided theorem proving systems, evaluation benchmarks, and hybrid architectures.

\subsubsection{LLM-Guided Theorem Proving Systems}

DeepSeek-Prover-V2 (2025) achieves state-of-the-art 88.9\% success on miniF2F-test through large-scale pretraining on formal mathematics corpora and Monte Carlo tree search. AlphaProof \citep{DeepMind2024} reaches IMO medal-level performance by combining informal problem statements with multi-step proof search. LeanCopilot integrates LLMs into the Lean proof assistant, enabling tactic suggestions from natural language goals. These systems demonstrate LLM effectiveness but do not isolate which mechanisms drive their advantage.

Our NL ablation provides mechanistic evidence for why these systems succeed: the 29.5 percentage point contribution from NL hints explains AlphaProof's reliance on informal problem statements and LeanCopilot's tactic suggestion accuracy. Where prior work optimizes architectures to maximize success rates, we decompose the advantage into testable mechanisms.

\subsubsection{Theorem Proving Benchmarks}

MiniF2F \citep{MiniF2F_2021} provides 488 Olympiad-level problems across Lean, Isabelle, and Metamath, establishing a standard evaluation benchmark. These benchmarks report LLM success rates but lack baselines for pure automated provers, making mechanistic comparisons impossible.

Our contribution: we establish the first measured lean-auto baseline on miniF2F Lean 4 subset (15.6\% [11.5\%, 20.3\%], N=244) and attribute ~60\% of the LLM gap to NL understanding. This baseline enables future controlled studies of architectural innovations.

\subsubsection{Hybrid LLM+Prover Systems}

Thor \citep{Thor_2022} combines LLMs with automated hammers, reporting 57.0\% total success with 8.2\% unique hybrid solutions (neither LLM nor hammer alone). This synergy suggests complementary strengths but does not explain when each component succeeds. Baldur \citep{Baldur_2023} uses LLMs to repair hammer-generated proof sketches. COPRA \citep{COPRA_2023} learns proof repair in Coq.

Our mechanistic attribution provides a routing hypothesis for hybrid systems: LLMs exploit NL hints (60\% advantage) while automated provers provide deterministic formal reasoning. Thor's 8.2\% unique solutions likely arise from problems with rich NL context (LLM strength) requiring reliable premise selection (hammer strength).

\subsubsection{Mechanistic Interpretability in Formal Methods}

Prior work in neural theorem proving focuses on improving LLM performance rather than explaining it. Polu et al. (2022) scale expert iteration on Lean, Lample et al. (2022) train Transformer models on Metamath, Mikula et al. (2023) apply retrieval augmentation---all report success metrics without mechanistic decomposition.

Our work is the first to apply mechanistic interpretability to theorem proving evaluation, isolating NL understanding (29.5 percentage point contribution, p<$10^{-9})$ and validating tactic budget as a fairness control metric (CV=0.36).
